\documentclass[11pt]{article}
\usepackage[margin=0.9in]{geometry}
\usepackage{amsmath,amssymb,amsthm,booktabs,array}
\usepackage[T1]{fontenc}
\usepackage{lmodern}
\usepackage[hidelinks]{hyperref}
\newcommand{\Ev}{\operatorname{Ev}}
\newcommand{\CT}{\operatorname{CT}}
\newcommand{\R}{\mathcal R}
\newtheorem{theorem}{Theorem}
\newtheorem{lemma}[theorem]{Lemma}
\newtheorem{corollary}[theorem]{Corollary}
\setlength{\parindent}{0pt}
\setlength{\parskip}{5pt}
\title{Weighted and rank-two extensions of Patel's proof: Hemmer--Straub--Westrem's $\Ev$ character-average conjecture}
\author{Haobo Ma\\\small ChronoAI Pte Ltd, Singapore
\and Wenlin Zhang\\\small National University of Singapore\\\small ChronoAI Pte Ltd, Singapore}
\date{16 November 2026}
\begin{document}
\maketitle
\begin{abstract}
Extending Patel's pointwise proof (2026, Solution 1), we give an alternative
coefficient proof of Conjecture 17, equation (6.1), page 15, of Hemmer,
Straub and Westrem's
\emph{New identities in the character table of symmetric groups involving
Riordan numbers} (2026), recovering the two-binary-quadratic Hilbert series
of Benanti et al. (2012, Example 3.3, p.~160).
Both weighted character sums vanish in odd degree and equal
$\binom{n/2+2}{2}$ in even degree.
Patel's three-variable Schur specialization first identifies the sums for
each partition, retaining all multiplicities in $\Ev$.
The cycle-index generating function then reduces their weighted average
to a constant term equal to $(1-t^2)^{-3}$.
We give explicit arbitrary-weight and common-divisor consequences, including
the $N=1$ case of Amdeberhan's $q$-series conjecture, and rederive the
Riordan tableau formula already implicit in Meilhan--Suzuki (2015).
A companion adjoint-representation argument proves the $N=2$ instances
of the source's two character identities and its $q$-series conjecture.
\end{abstract}

\section{The weighted identity}
Hemmer, Straub and Westrem~\cite{HSW} ask for an explicit evaluation of
two centralizer-weighted sums of symmetric-group characters.
Their Conjecture 17 is Conjecture 6.1 on page 12 of the arXiv version;
equation (6.1) has the same numbering in both versions.
We evaluate the common sum by first reducing each summand to adjacent
coefficients of a Laurent polynomial. This reduction also gives their
pointwise Conjecture 6, equation (2.4) (arXiv Conjecture 2.3).

For $\lambda=(a_1,\ldots,a_r)\vdash n$, let $m_k$ be the multiplicity
of $k$ and set $z_\lambda=\prod_{k\geq1} k^{m_k}m_k!$.
The multiset $\Ev(\lambda)$ is obtained by independently replacing each
occurrence of $a_i$ with either $(2a_i)$ or $(a_i,a_i)$, and sorting.
All $2^r$ choices are counted, even when they give identical partitions.
Write $\chi^\mu_\tau$ for the irreducible character of $S_{2n}$ indexed
by $\mu$, evaluated at cycle type $\tau$.
Let $\R_3(2n)$ consist of partitions of $2n$ with at most three parts,
all even; zero padding to a triple is understood.
Define
\begin{align*}
 L_\lambda&=\sum_{\tau\in\Ev(\lambda)}(-1)^{\ell(\tau)}
                         \sum_{\mu\in\R_3(2n)}\chi^\mu_\tau,&
 B_\lambda&=\sum_{\tau\in\Ev(\lambda)}\chi^{(n,n)}_\tau,\\
 L_n&=\sum_{\lambda\vdash n}\frac{L_\lambda}{z_\lambda},&
 B_n&=\sum_{\lambda\vdash n}\frac{B_\lambda}{z_\lambda}.
\end{align*}

\begin{theorem}\label{thm:weighted}
For every integer $n\geq1$,
\[
 L_n=B_n=
 \begin{cases}\displaystyle\binom{n/2+2}{2},&n\text{ even},\\
 0,&n\text{ odd}.
 \end{cases}
\]
With the empty-partition convention $L_0=B_0=1$, the generating function is
\begin{equation}\label{eq:answer}
 \sum_{n\geq0}L_nt^n=\sum_{n\geq0}B_nt^n=(1-t^2)^{-3}.
\end{equation}
\end{theorem}

Patel's public solution~\cite[Solution 1, Sections 1--3]{P} proves the full
pointwise identity, including the Schur detector and the same constant-term
formula used here. We retain its complete coefficient argument and develop
the weighted consequences. The cycle-index extraction is also the Hilbert
series of two binary quadratics, evaluated by Benanti et al.
\cite[Example 3.3, p.~160]{BDG}; the precise identification is given below.
Thus the weighted evaluation is an application of known invariant theory,
not a first evaluation of that Hilbert series.

The posing paper proves the pointwise identity for partitions with a
single part size, obtaining Riordan numbers or central trinomial
coefficients~\cite[Sections 4--5]{HSW}; we rederive these specializations below.
Westrem's signed power-sum factorization~\cite[Theorem 2.1]{W}
is a starting ingredient. A representation-theoretic antecedent of the
same-parity detector is Meilhan and Suzuki's decomposition of invariant
tensor powers~\cite[Lemma 3.7]{MS}. The later character-degree
refinements of Hemmer, Straub and Westrem~\cite[Theorem 1.3]{K}
concern a different specialization, rather than these weighted cycle-type sums.
The related Riordan tableau enumeration follows from that decomposition
and is also contained in Okada~\cite[Corollary 5.5(3)]{O}, with a later bijective proof by
Jagenteufel~\cite{J}. Amaral, Gil and Weiner~\cite[Theorem 4.1]{AGW}
subsequently answered the source's full Motzkin-level bijection question.

\section{The pointwise reduction}
Brackets denote coefficient extraction, and $\CT_u=[u^0]$ denotes the
constant coefficient of a Laurent polynomial in $u$.
Put $p_a=x^a+y^a+z^a$ and $\Delta=(x-y)(x-z)(y-z)$.
Let $\mathcal S_3(m)$ be the weakly decreasing triples of nonnegative
integers summing to $m$ whose entries have the same parity.
Thus $\mathcal S_3(2n)=\R_3(2n)$, while $\mathcal S_3(m)$ has all entries
odd when $m$ is odd. For a symmetric homogeneous polynomial $F$ of degree
$m\geq0$, define
\[
 E_m(F)=\sum_{\mu\in\mathcal S_3(m)}
 [x^{\mu_1+2}y^{\mu_2+1}z^{\mu_3}]\,\Delta F.
\]

\begin{lemma}[Patel's Schur detector]\label{lem:detector}
For every symmetric homogeneous $F\in\mathbb Q[x,y,z]$ of degree $m\geq0$,
\begin{equation}\label{eq:detector}
 E_m(F)=([u^m]-[u^{m+1}])F(u^2,u,1).
\end{equation}
\end{lemma}
\begin{proof}
Expand $F$ in the Schur basis $s_{(a,b,c)}$, where
$a\geq b\geq c\geq0$ and $a+b+c=m$.
The alternant formula
$\Delta s_{(a,b,c)}=\det(x_i^{\mu_j+3-j})_{1\leq i,j\leq3}$,
with $(x_1,x_2,x_3)=(x,y,z)$ and $\mu=(a,b,c)$,
shows that $E_m(s_{(a,b,c)})$ is $1$ if all three parts have the same parity and $0$
otherwise: the exponent triple $(a+2,b+1,c)$ is strictly decreasing,
and no other permutation of it is a decreasing triple.

Set $d=a-b$ and $e=b-c$.
Evaluating the same alternant at $(u^2,u,1)$ gives
\[
 H(u):=s_{(a,b,c)}(u^2,u,1)
 =u^{b+2c}
 \frac{(1-u^{d+1})(1-u^{e+1})(1-u^{d+e+2})}
 {(1-u)^2(1-u^2)}.
\]
This equality may be read in $\mathbb Q[[u]]$, since its denominator
has constant coefficient $1$.
Symmetry and homogeneity imply $H(u)=u^{2m}H(u^{-1})$, so
$[u^{m+1}]H=[u^{m-1}]H$.
Since $m-(b+2c)=d+e$, the right side of~\eqref{eq:detector} on this basis element is
\[
 [u^{d+e}]\frac{(1-u^{d+1})(1-u^{e+1})(1-u^{d+e+2})}
 {(1-u)(1-u^2)}.
\]
Let $K(j)=\lfloor j/2\rfloor+1$ for $j\geq0$, and $K(j)=0$ for $j<0$.
The last numerator factor and the product $u^{d+1}u^{e+1}$ have degree
greater than $d+e$; the displayed coefficient therefore equals
\[
 K(d+e)-K(e-1)-K(d-1)
 =\begin{cases}1,&d,e\text{ both even},\\0,&\text{otherwise}.\end{cases}
\]
Indeed, substituting $d=2i+\epsilon$, $e=2j+\eta$ reduces this to the
four choices $\epsilon,\eta\in\{0,1\}$.
Both gaps are even exactly when $a,b,c$ have the same parity.
Thus the two functionals agree on every Schur basis element.
\end{proof}

\begin{lemma}[Patel's pointwise identity]\label{lem:pointwise}
For every $\lambda=(a_1,\ldots,a_r)\vdash n$ with $n\geq1$, put
$Q_\lambda(u)=\prod_{i=1}^r(u^{-a_i}+1+u^{a_i})$. Then
\begin{equation}\label{eq:pointwise}
 L_\lambda=B_\lambda=2^r\CT_u\bigl((1-u^{-1})Q_\lambda(u)\bigr).
\end{equation}
\end{lemma}
\begin{proof}
Define $F_- =\prod_i(p_{a_i}^2-p_{2a_i})$ and
$F_+ =\prod_i(p_{a_i}^2+p_{2a_i})$.
Expanding the choices defining $\Ev$ gives
$F_-=\sum_{\tau\in\Ev(\lambda)}(-1)^{\ell(\tau)}p_\tau$ and
$F_+=\sum_{\tau\in\Ev(\lambda)}p_\tau$.
The first is Westrem's factorization~\cite{W} in three variables;
the expansion also makes the multiplicities explicit.
The Frobenius character formula~\cite[Lemma 7 and equation (5.2)]{HSW}
identifies
\[
 \chi^\mu_\tau=[x^{\mu_1+2}y^{\mu_2+1}z^{\mu_3}]\,\Delta p_\tau
 \qquad(\ell(\mu)\leq3).
\]
Consequently $L_\lambda=E_{2n}(F_-)$ and
$B_\lambda=[x^{n+2}y^{n+1}z^0]\Delta F_+$.
Let $P_\lambda(u)=\prod_i(1+u^{a_i}+u^{2a_i})=u^nQ_\lambda(u)$.
Direct expansion of a single factor yields
\[
 F_-(u^2,u,1)=2^ru^nP_\lambda(u),\qquad
 F_+(u,1,0)=2^rP_\lambda(u).
\]
Lemma~\ref{lem:detector} now gives
$L_\lambda=2^r([u^n]-[u^{n+1}])P_\lambda$.
For $B_\lambda$, only the $z$-free part $x^2y-xy^2$ of $\Delta$
contributes. Since $n\geq1$, homogeneity gives
\begin{align*}
 B_\lambda
 &=([x^ny^n]-[x^{n+1}y^{n-1}])F_+(x,y,0)\\
 &=2^r([u^n]-[u^{n+1}])P_\lambda(u).
\end{align*}
Finally $P_\lambda=u^nQ_\lambda$, so the adjacent-coefficient difference
is $[u^0]Q_\lambda-[u^1]Q_\lambda$
$=\CT_u((1-u^{-1})Q_\lambda)$, as required.
\end{proof}

For example, $\lambda=(2,1)$ gives
$P_\lambda=1+u+2u^2+u^3+2u^4+u^5+u^6$ and
$L_\lambda=B_\lambda=4(1-2)=-4$.
Thus the nonnegative weighted averages in Theorem~\ref{thm:weighted} do not
reflect positivity of every individual summand.

\section{Short consequences and earlier special cases}
Let $R_m$ count Riordan paths of length $m$: paths from $(0,0)$ to
$(m,0)$ with steps $(1,1),(1,0),(1,-1)$, never below the horizontal
axis, with no horizontal step on that axis. Let
$T_r=[u^0](u^{-1}+1+u)^r$ be the central trinomial coefficient.
For a partition $\mu$, let $f^\mu$ be the number of standard Young tableaux
of shape $\mu$, equivalently the degree of $\chi^\mu$.
The adjacent-trinomial description of $R_m$ below is Callan's formula
\cite{C}, recorded earlier in his OEIS comment dated 7 February 2004~\cite{OEIS}.

\begin{corollary}[Riordan tableaux and equal parts]\label{cor:riordan}
For $m\geq0$ and positive integers $k,r$,
\[
 \sum_{\mu\in\mathcal S_3(m)} f^\mu=R_m,\qquad
 L_{(k^r)}=B_{(k^r)}=
 \begin{cases}2^rR_r,&k=1,\\2^rT_r,&k>1.\end{cases}
\]
The first formula follows already from Meilhan--Suzuki
\cite[Lemma 3.7]{MS} and Okada~\cite[Corollary 5.5(3)]{O}.
For the second, the rectangular $k=1$ case is recorded in Wallach's OEIS
comment dated 20 August 2014~\cite{OEIS}; both $\Ev$ sums and the $k>1$
case are proved in~\cite[Sections 4--5]{HSW}.
\end{corollary}
\begin{proof}
Apply Lemma~\ref{lem:detector} to $F=(x+y+z)^m$ and use the
Frobenius formula at the identity class. This gives
$\sum_{\mu\in\mathcal S_3(m)}f^\mu
=\CT_u(1-u^{-1})(u^{-1}+1+u)^m$.
To identify the latter with $R_m$, write
$D(s)=\sqrt{1-2s-3s^2}$, with constant term $1$.
Expanding in $s$ gives
\[
 A(s):=\CT_u\frac1{1-s(1+u+u^{-1})}=
 \frac1{1-s}\sum_{j\geq0}\binom{2j}{j}
       \left(\frac{s}{1-s}\right)^{2j}=D(s)^{-1}.
\]
Multiplying the denominator through and using symmetry in $u,u^{-1}$
shows that the adjacent constant term is
$((1-s)A(s)-1)/(2s)$. Their difference is therefore
\[
 \CT_u\frac{1-u^{-1}}{1-s(1+u+u^{-1})}
 =\frac{1+(3s-1)/D(s)}{2s}
 =\frac{1+s-D(s)}{2s(1+s)}.
\]
For completeness, arbitrary Motzkin excursions have series
$M=1+sM+s^2M^2$, so $s^2M=(1-s-D(s))/2$.
First-return decomposition of a nonempty Riordan path gives
$R=1+s^2MR$, hence
$R=2/(1+s+D(s))=(1+s-D(s))/(2s(1+s))$.
All displayed quotients have their indicated formal-series expansion at
$s=0$. Finally, apply~\eqref{eq:pointwise} to $(k^r)$.
For $k=1$ the adjacent-coefficient difference is $R_r$; for $k>1$
every exponent is divisible by $k$, so the coefficient of $u$ is zero
and the constant coefficient is $T_r$.
\end{proof}

\begin{corollary}[Common divisors]\label{cor:gcd}
If $\gcd(a_1,\ldots,a_r)>1$, then
\[
 L_\lambda=B_\lambda
 =2^r\#\{(\epsilon_1,\ldots,\epsilon_r)\in\{-1,0,1\}^r:
                   \textstyle\sum_i\epsilon_i a_i=0\}\geq0.
\]
This extends the central-trinomial case to mixed part sizes.
\end{corollary}
\begin{proof}
Every exponent of $Q_\lambda$ is divisible by the common divisor,
so $[u^1]Q_\lambda=0$. Its constant coefficient counts the displayed
choices. Apply Lemma~\ref{lem:pointwise}.
\end{proof}

\begin{corollary}[Vanishing above the first-row bound]\label{cor:row}
For $\lambda\vdash n$ and $\mu\vdash2n$ with $\mu_1>n$,
\[
 \sum_{\tau\in\Ev(\lambda)}(-1)^{\ell(\tau)}\chi^\mu_\tau=0.
\]
This is the row-wise vanishing consequence of Westrem's
factorization~\cite[Corollary 3.4]{W}.
\end{corollary}
\begin{proof}
In $d\geq\ell(\mu)$ variables the signed power-sum sum is
$F_-=\prod_i2\sum_{j<k}x_j^{a_i}x_k^{a_i}$.
Each variable has degree at most $n$. Multiplying by the Vandermonde
adds at most $d-1$ to any variable's degree, whereas the Frobenius
coefficient for $\mu$ requires exponent $\mu_1+d-1>n+d-1$
on the first variable. That coefficient is zero.
\end{proof}

\section{Cycle-index averaging}
\begin{proof}[Proof of Theorem~\ref{thm:weighted}]
Work in $\mathbb Q[u,u^{-1}][[t]]$.
All constant terms below are taken coefficientwise in $t$.
For $a_k(u)=2(u^k+1+u^{-k})$, define
\[
 C(t,u)=\sum_{n\geq0}t^n\sum_{\lambda\vdash n}
                  \frac{\prod_i a_{\lambda_i}(u)}{z_\lambda}.
\]
Writing a partition in terms of its multiplicities gives the formal identity
\begin{align}
 C(t,u)
 &=\prod_{k\geq1}\sum_{j\geq0}
       \frac{(a_k(u)t^k)^j}{k^j j!}
 =\exp\!\left(\sum_{k\geq1}\frac{a_k(u)t^k}{k}\right)\notag\\
 &=\frac{1}{(1-t)^2(1-tu)^2(1-t/u)^2}.\label{eq:cycle}
\end{align}
Every coefficient of $t^n$ involves only $k\leq n$ and finitely many
multiplicities, so the products, sums and constant-term operations are
well-defined. The last equality follows from
$\exp(2\sum_{k\geq1}v^k/k)=(1-v)^{-2}$ in formal power series.
Lemma~\ref{lem:pointwise}, with the empty partition handled separately,
therefore gives
\[
 \sum_{n\geq0}L_nt^n=\sum_{n\geq0}B_nt^n
 =\frac{1}{(1-t)^2}\CT_u
       \frac{1-u^{-1}}{(1-tu)^2(1-t/u)^2}.
\]
Expand the last two factors using
$(1-v)^{-2}=\sum_{j\geq0}(j+1)v^j$.
The unshifted constant term requires equal exponents; the shifted term
requires the exponent from $(1-tu)^{-2}$ to be one larger. Hence
\begin{align*}
 \CT_u\frac{1-u^{-1}}{(1-tu)^2(1-t/u)^2}
 &=\sum_{j\geq0}(j+1)^2t^{2j}
       -\sum_{j\geq0}(j+1)(j+2)t^{2j+1}\\
 &=\frac{1+t^2}{(1-t^2)^3}-\frac{2t}{(1-t^2)^3}
 =\frac{(1-t)^2}{(1-t^2)^3}.
\end{align*}
The two rational-series formulas follow by differentiating
$\sum_{j\geq0}v^j=(1-v)^{-1}$.
Cancelling $(1-t)^2$ proves~\eqref{eq:answer}.
Finally $(1-t^2)^{-3}=\sum_{j\geq0}\binom{j+2}{2}t^{2j}$,
which gives both the even value and the vanishing of odd coefficients.
\end{proof}

\subsection*{The published invariant Hilbert series}
Let $V=\operatorname{Sym}^2\mathbb C^2$ be the binary-quadratic
representation of $SL_2(\mathbb C)$, with torus weights $u,1,u^{-1}$
(writing $u$ for the square of the usual torus parameter).
The graded character of the symmetric algebra on $V\oplus V$ is
\[
 \sum_{n\geq0}\operatorname{ch}\operatorname{Sym}^n(V\oplus V)t^n
 =C(t,u).
\]
Every irreducible constituent factors through $SO_3$ and has character
$u^{-j}+\cdots+u^j$ for some integer $j\geq0$.
The functional $\CT_u(1-u^{-1})$ takes value $1$ for $j=0$ and $0$
otherwise, hence extracts the dimension of the invariant subspace.
Consequently
\[
 \sum_{n\geq0}B_nt^n
 =H\!\left(\mathbb C[V\oplus V]^{SL_2};t\right)
 =(1-t^2)^{-3}.
\]
Here $H(A;t)=\sum_n\dim(A_n)t^n$ denotes the Hilbert series; coordinate
functions have the same character since $V$ is self-dual.
The last series is exactly Benanti et al.'s evaluation
\cite[Example 3.3, p.~160]{BDG}, with their $W(2)=V$.
The elementary calculation above supplies an alternative coefficient proof
and its explicit interpretation as the weighted $\Ev$ character sums.

\section{Arbitrary weights and the source's \texorpdfstring{$q$}{q}-series}
\begin{corollary}\label{cor:q}
With $g_\lambda(q)=\prod_i q^{a_i}/(1+q^{a_i})$,
\begin{equation}\label{eq:q}
 \sum_{n\geq0}\sum_{\lambda\vdash n}
       \frac{g_\lambda(q)}{z_\lambda}L_\lambda
 =\sum_{n\geq0}\sum_{\lambda\vdash n}
       \frac{g_\lambda(q)}{z_\lambda}B_\lambda.
\end{equation}
Here $L_\varnothing=B_\varnothing=1$.
This is the $N=1$ case of Amdeberhan's Conjecture 18,
equation (6.2), page 16 of~\cite{HSW} (arXiv Conjecture 6.2).
\end{corollary}
\begin{proof}
For $n\geq1$, multiply~\eqref{eq:pointwise} by
$g_\lambda(q)/z_\lambda$ and sum.
Since $g_\lambda(q)$ has order at least $|\lambda|$, each coefficient
of $q$ receives contributions from only finitely many partitions.
The empty-partition terms are $1$ on both sides.
\end{proof}

\begin{corollary}[Arbitrary multiplicative weights]\label{cor:weights}
Let $w_k$ be elements of a commutative $\mathbb Q$-algebra, with $t$ a
formal variable. The two series
$\sum_{n\geq0}t^n\sum_{\lambda\vdash n}z_\lambda^{-1}
 (\prod_iw_{a_i})L_\lambda$ and the corresponding series with $B_\lambda$
are equal, and their common value is
\[
 \CT_u(1-u^{-1})\exp\!\left(
 2\sum_{k\geq1}\frac{w_k t^k(u^k+1+u^{-k})}{k}\right).
\]
\end{corollary}
\begin{proof}
Multiply~\eqref{eq:pointwise} by the indicated weights and sum by the
multiplicities of each part, exactly as in~\eqref{eq:cycle}.
Each coefficient of $t$ involves finitely many partitions and Laurent
monomials, so both operations are defined coefficientwise.
\end{proof}

For the specific $q$-weights, multiplicity summation also gives the common
series in~\eqref{eq:q} as
\[
 \CT_u(1-u^{-1})\exp\!\left(
 2\sum_{k\geq1}\frac{q^k(u^k+1+u^{-k})}{k(1+q^k)}\right).
\]
Its coefficients through degree $10$ are
\[
 1+3q^2-4q^3+9q^4-12q^5+22q^6-36q^7+60q^8-88q^9+135q^{10},
\]
agreeing with the finite expansion in~\cite[p.~16]{HSW}.

The enumeration in Corollary~\ref{cor:riordan} follows by taking dimensions
in Meilhan--Suzuki's invariant-tensor decomposition~\cite[Lemma 3.7]{MS}
(2015). It also follows from Okada~\cite[Corollary 5.5(3)]{O} (2016):
specialize the rank to $1$ and the
terminal column size to $0$. His partition sequences become Riordan
paths, and the permitted shapes have zero or three odd rows.
Jagenteufel~\cite[Section 3.2.1 and Lemma 36]{J} later supplied an explicit
Riordan-path/tableau bijection. Hemmer, Straub and Westrem subsequently
gave an independent enumerative proof~\cite[Theorem 11]{HSW}.
The stronger question on page 16 of~\cite{HSW}, repeated in~\cite{K},
asks for a full Motzkin-path/tableau bijection restricting to these Riordan
objects. Amaral, Gil and Weiner~\cite[Theorems 2.1, 3.7 and 4.1]{AGW}
answer it affirmatively, using restricted permutations and a tableau
parity correction. These are prior results, not constructions of this note.

Questions 4 and 5 are classification questions, not universal assertions.
Their unrestricted equalities already fail in the posing source: equation
(2.1) has values $1040$ and $1041$ for $n=12,N=3$; Appendix A gives a
pointwise discrepancy of $8$ for $\lambda=(5,2,1),N=3$.
Our answers for $N=1,2$ do not change those counterexamples.

\subsection*{The rank-two companion}
Write $\mathcal R_d(2n)$ for the even-row partitions of $2n$ with at most
$d$ rows, and $\mathcal R_d^c(2n)$ for the partitions with at most $d$
rows and all columns even.
For a virtual character $f$ of a reductive group $G$, write
$\langle f\rangle_G$ for its trivial multiplicity, extended linearly
from actual representations.

\begin{corollary}\label{cor:ranktwo}
For every $\lambda\vdash n$, $n\geq1$,
\[
 \sum_{\tau\in\Ev(\lambda)}(-1)^{\ell(\tau)}
        \sum_{\mu\in\mathcal R_5(2n)}\chi^\mu_\tau
 =\sum_{\tau\in\Ev(\lambda)}
        \sum_{\mu\in\mathcal R_4^c(2n)}\chi^\mu_\tau.
\]
Thus Questions 4 and 5 have affirmative answers also for $N=2$, and
Conjecture 18 holds for $N=2$ as an equality of formal $q$-series.
\end{corollary}
\begin{proof}
Let $X=\mathbb C^5$ have a nondegenerate symmetric form and
$W=\mathbb C^4$ a nondegenerate alternating form.
The orthogonal and symplectic invariant criteria
\cite[Corollaries 3.3 and 2.3]{DH} give
$\langle s_\mu(X)\rangle_{O_5}=1$ exactly for even rows, and
$\langle s_\mu(W)\rangle_{Sp_4}=1$ exactly for even columns.
Both multiplicities are zero otherwise, including the dimension bounds.
Apply these criteria to the Frobenius expansions of $F_-$ in five
variables and $F_+$ in four variables.
For any representation $U$, define its Adams character by
$\psi^a(U)(g)=\operatorname{tr}(g^a|U)$; these are virtual characters.
The power-sum identities
\[
 p_a(X)^2-p_{2a}(X)=2\psi^a(\Lambda^2X),\qquad
 p_a(W)^2+p_{2a}(W)=2\psi^a(\operatorname{Sym}^2W)
\]
show that the two sums are respectively
$2^r\langle\prod_i\psi^{a_i}(\Lambda^2X)\rangle_{O_5}$ and
$2^r\langle\prod_i\psi^{a_i}(\operatorname{Sym}^2W)\rangle_{Sp_4}$.
The action of $O_5$ on $\Lambda^2X$ has the same image as that of
$SO_5$: if $\det g=-1$, then $-g\in SO_5$ and
$\Lambda^2g=\Lambda^2(-g)$.

Here is the exceptional identification needed to compare these characters.
Let $b\in\Lambda^2W$ be the inverse symplectic bivector. The wedge pairing
on $\Lambda^2W$ is symmetric and nondegenerate, and $b\wedge b\ne0$.
Its orthogonal complement $V=b^\perp$ is therefore five-dimensional
and nondegenerate. The action induces $\rho:Sp_4\to SO(V)$.
Identify $X$ with $V$ by an isometry.
If an infinitesimal symplectic endomorphism acts trivially on $V$, it
also fixes $b$, and hence acts trivially on $\Lambda^2W$.
The latter forces it to be zero: in a basis, all off-diagonal entries
vanish and its diagonal entries satisfy $d_i+d_j=0$ for every $i\ne j$.
Thus $d\rho$ is injective; both Lie algebras have dimension $10$,
so it is an isomorphism and $\rho$ is onto the connected group $SO(V)$.
The natural equivariant identifications
$\Lambda^2V\simeq\mathfrak{so}(V)$ and
$\operatorname{Sym}^2W\simeq\mathfrak{sp}(W)$, together with $d\rho$,
identify the two adjoint representations after pullback.
Their Adams characters and products agree, and surjectivity preserves
trivial multiplicities. This proves the displayed identity.
Multiplication by $1/z_\lambda$ and summation proves Question 4;
multiplication by $g_\lambda(q)/z_\lambda$ proves Conjecture 18,
with coefficientwise finiteness as in Corollary~\ref{cor:q}.
\end{proof}

\begin{center}
\small
\begin{tabular}{@{}>{\raggedright\arraybackslash}p{0.39\linewidth}>{\raggedright\arraybackslash}p{0.54\linewidth}@{}}
\toprule
Source item in~\cite{HSW} & Coverage\\
\midrule
Conjecture 17, (6.1), p.~15 & Alternative proof; prior equivalent Hilbert series Benanti et al. (2012); Theorem~\ref{thm:weighted}.\\
Conjecture 6, (2.4), p.~4 & Reproved; prior proof Patel (2026); Lemma~\ref{lem:pointwise} includes all mixed part sizes.\\
Question 4, (2.1), pp.~3--4 & Answered: yes for $N=1,2$; universally no, by the source's $N=3$ counterexample.\\
Question 5, (2.3), p.~4 & Answered: yes for $N=1,2$; universally no, by the source's Appendix A.\\
Riordan tableau formula, Theorem 11 & Reproved; prior proof Meilhan--Suzuki (2015); also Okada (2016), Jagenteufel (2018), Hemmer--Straub--Westrem (2025).\\
Single-part-size cases, Sections 4--5 & Reproved; prior rectangular $k=1$ formula Wallach (2014); both sides Hemmer--Straub--Westrem (2025).\\
Conjecture 18, (6.2), p.~16 & Proved for $N=1,2$ by Corollaries~\ref{cor:q} and~\ref{cor:ranktwo}; no assertion for $N\geq3$.\\
Full Motzkin--Riordan bijection question, p.~16 & Answered: yes; prior proof Amaral--Gil--Weiner (2026), Theorem 4.1.\\
Westrem, Corollary 3.4 & Reproved; prior proof Westrem (2024); Corollary~\ref{cor:row}.\\
Mixed common-divisor case; arbitrary weights & Not posed; proved in Corollaries~\ref{cor:gcd} and~\ref{cor:weights}.\\
Other $N\geq3$ classifications and degree refinements of~\cite{K} & Not covered by the present argument.\\
\bottomrule
\end{tabular}
\end{center}

\section{Mathematical verification scope}
At the pinned sources,
\href{https://github.com/the-omega-institute/trureturing/blob/439c57ca26fd0c0d2de1f48bc492717bc6f7597c/D5/S3/Combinatorics/Partitions/HemmerStraubWestremWeightedAverage.lean}{\texttt{HemmerStraubWestremWeightedAverage.lean}}
contains the declaration \texttt{result} proving both weighted equalities
for every $n\geq1$ in its rational Frobenius-coefficient model.
The declaration \texttt{result} in
\href{https://github.com/the-omega-institute/trureturing/blob/439c57ca26fd0c0d2de1f48bc492717bc6f7597c/D5/S3/Combinatorics/Partitions/HemmerStraubWestrem.lean}{\texttt{HemmerStraubWestrem.lean}}
proves the pointwise identity for every positive-size partition.
In these files, \texttt{frobChar} is the coefficient of
$\Delta p_\tau$ used above. The interpretation as ordinary irreducible
characters is supplied by the classical Frobenius formula, rather than an
internal construction of Specht modules.
The all-degree Schur-basis argument of Lemma~\ref{lem:detector}, its
same-parity extension to odd degree, and the corollaries on Riordan
tableaux, common divisors, first-row vanishing, arbitrary weights and
$q$-series are written proofs in this note. The cited path bijections
have the scopes stated above.
Corollary~\ref{cor:ranktwo}, including its $N=2$ weighted and $q$-series
consequences, is a written representation-theoretic proof.

The accompanying \texttt{verify.py} uses Murnaghan--Nakayama recursion to
check both character sums and their Laurent formula for $0\leq n\leq8$.
It also compares every Laurent coefficient in the cycle recurrence and
the product~\eqref{eq:cycle} through degree $60$, and checks the displayed
$q$-series through degree $10$. It checks the additional corollaries
over the finite ranges specified in the accompanying README.
All computations use exact integers and
rational numbers; the all-degree conclusions are established by the proofs.

\begin{thebibliography}{99}
\bibitem{DH}
V. Drensky and E. Hristova, \emph{Invariants of symplectic and orthogonal
groups acting on $GL(n,\mathbb C)$-modules},
\href{https://arxiv.org/abs/1707.05893v1}{arXiv:1707.05893v1}
(18 July 2017), Corollaries 2.3 and 3.3.
\bibitem{P}
S. Patel, \emph{A proof of the character-sum identity for every partition},
\href{https://mathdb.com/p/370115/the-n-1-conjecture-for-character-sums-over-riordan-partition}{MathDB problem 370115},
Solution 1, Sections 1--3, public posting (2026), accessed 16 November 2026.
\bibitem{BDG}
F. Benanti, S. Boumova, V. Drensky, G. K. Genov and P. Koev,
\emph{Computing with rational symmetric functions and applications to
invariant theory and PI-algebras}, Serdica Mathematical Journal
\textbf{38} (2012), 137--188, Example 3.3, p.~160.
\href{https://arxiv.org/abs/1201.4448v1}{arXiv:1201.4448v1}
(21 January 2012).
\bibitem{MS}
J.-B. Meilhan and S. Suzuki, \emph{Riordan trees and the homotopy
$sl_2$ weight system}, \href{https://arxiv.org/abs/1507.04454v1}{arXiv:1507.04454v1}
(16 July 2015), Section 3.3, Lemma 3.7; Journal of Pure and Applied Algebra
\textbf{221}(3) (2017), 691--706,
\href{https://doi.org/10.1016/j.jpaa.2016.07.012}{doi:10.1016/j.jpaa.2016.07.012}.
\bibitem{C}
D. Callan, \emph{Riordan numbers are differences of trinomial coefficients},
\href{https://pages.stat.wisc.edu/~callan/notes/riordan/riordan.pdf}{author manuscript}
(25 September 2006), 2 pp.
\bibitem{OEIS}
OEIS Foundation, \href{https://oeis.org/A005043}{A005043, Riordan numbers},
comments attributed to D. Callan (7 February 2004) and N. Wallach
(20 August 2014), accessed 16 November 2026.
\bibitem{O}
S. Okada, \emph{Pieri rules for classical groups and equinumeration between
generalized oscillating tableaux and semi-standard tableaux},
Electronic Journal of Combinatorics \textbf{23}(4) (2016), P4.43,
\href{https://doi.org/10.37236/6214}{doi:10.37236/6214}.
Preprint: \href{https://arxiv.org/abs/1606.02375v1}{arXiv:1606.02375v1},
Corollary 5.5(3).
\bibitem{AGW}
R. J. Amaral, J. B. Gil and M. D. Weiner,
\emph{Motzkin paths, 321-avoiding permutations, and standard Young tableaux
with rows of equal parity},
\href{https://arxiv.org/abs/2609.27096v1}{arXiv:2609.27096v1}
(22 September 2026), Theorem 4.1;
\href{https://arxiv.org/abs/2609.27096v2}{version 2} (28 September 2026).
\bibitem{HSW}
D.~J. Hemmer, A. Straub and K.~J. Westrem,
\emph{New identities in the character table of symmetric groups involving
Riordan numbers}, Electronic Journal of Combinatorics \textbf{33}(1)
(2026), P1.57, 17 pp.,
\href{https://doi.org/10.37236/14401}{doi:10.37236/14401}.
Preprint: \href{https://arxiv.org/abs/2509.02796v1}{arXiv:2509.02796v1} (2025).
\bibitem{W}
K.~J. Westrem, \emph{A new symmetric function identity with an application
to symmetric group character values},
\href{https://arxiv.org/abs/2410.04644v1}{arXiv:2410.04644v1} (2024), Theorem 2.1.
\bibitem{K}
D.~J. Hemmer, A. Straub and K.~J. Westrem,
\emph{Equal knapsack identities between symmetric group character degrees},
\href{https://arxiv.org/abs/2510.00301v2}{arXiv:2510.00301v2} (2026).
\bibitem{J}
J. Jagenteufel, \emph{A Sundaram type bijection for $\mathrm{SO}(3)$:
vacillating tableaux and pairs of standard Young tableaux and orthogonal
Littlewood--Richardson tableaux}, Electronic Journal of Combinatorics
\textbf{25}(3) (2018), P3.50,
\href{https://doi.org/10.37236/7713}{doi:10.37236/7713}.
Preprint: \href{https://arxiv.org/abs/1801.03780v2}{arXiv:1801.03780v2},
Section 3.2.1 and Lemma 36.
\end{thebibliography}

\smallskip
\small\copyright\ 2026 Haobo Ma and Wenlin Zhang.
Licensed under \href{https://creativecommons.org/licenses/by/4.0/}{Creative Commons Attribution 4.0 International (CC BY 4.0)}.
\end{document}
